\section{Coordinates, Points, and Directions}

Every number shown in the Item panel is a coordinate, and a coordinate has meaning only relative to a frame. Blender uses a right-handed coordinate system in which $Z$ points up. The world origin is the point $(0,0,0)$, and the three world axes, drawn in red, green, and blue in the default theme, are the $X$, $Y$, and $Z$ axes. Right-handed means that if the fingers of the right hand curl from the positive $X$ axis toward the positive $Y$ axis, the thumb points along positive $Z$. A rotation about $Z$ by a positive angle therefore turns $X$ toward $Y$, which is the convention used in every rotation matrix in this book.

Two kinds of quantity are easily confused. A \emph{point} is a location, such as the object's origin or a vertex. A \emph{direction} or displacement is a quantity such as ``two units along $X$,'' which has length and orientation but no position. The Location field of an object is a point. Its Dimensions field is a set of extents, which behave like displacements. The difference becomes important when transformations are written as matrices, because translating a point moves it while translating a direction leaves it unchanged. Chapter 2 makes this precise.

Vectors can be added componentwise, scaled by a number, and measured. The length of $\mathbf v=(v_x,v_y,v_z)$ is $\|\mathbf v\|=\sqrt{v_x^2+v_y^2+v_z^2}$, and the distance between two points is the length of their difference. The dot product $\mathbf a\cdot\mathbf b=a_xb_x+a_yb_y+a_zb_z$ equals $\|\mathbf a\|\,\|\mathbf b\|\cos\theta$, so it is zero exactly when two nonzero vectors are perpendicular. The platform of Section~1.3 gives a concrete check. Its top face is a six by six square, so the diagonal across it has length $\sqrt{6^2+6^2}=\sqrt{72}\approx 8.485$ units, a value that can be confirmed by selecting two opposite top vertices and reading their coordinates.

Each object carries its own \emph{local} coordinate frame, with its origin at the point shown by the small orange dot when the object is selected. The vertices of the startup cube are stored in local coordinates as the eight points $(\pm1,\pm1,\pm1)$. Where those vertices appear in the scene, which is given by their \emph{world} coordinates, depends on the object's location, rotation, and scale. For an object without a parent, the Location field gives the world coordinates of the object's origin. In Edit Mode, the Item panel of the Sidebar reports the selected vertex or the median of the selected vertices and offers a choice between Local and Global display. Reading both for the same vertex is the quickest way to see that one point has different coordinates in different frames.

\section{Exercise: Reading the Platform as Numbers}

\begin{exercise}[1.6: Local and world coordinates]
Using the platform from Section~1.3, enter Edit Mode, select a single vertex, and open the Item panel. Record the vertex's Local and Global coordinates for at least three vertices, including one on the top face and one on the bottom face. Before looking, predict each Global coordinate from the object's scale $(3,3,0.2)$ and location $(0,0,-0.2)$, using the rule that each local coordinate is multiplied by the scale and the location is then added.

The vertex with local coordinates $(1,1,1)$ should have global coordinates $(3,3,0)$, and the vertex $(-1,-1,-1)$ should have global coordinates $(-3,-3,-0.4)$. State, in a sentence, why the top of the platform lies exactly at $Z=0$ and where the object's origin lies relative to the platform's thickness.
\end{exercise}
